\section{Senatorial Rome and a Household of Saints (c.~540--574)}

\subsection{A ruined peninsula}

Gregory grew up in the wreckage of imperial Italy. No contemporary source
records the exact year of his birth; modern scholarship places it around 540,
roughly a decade after Saint Benedict of Nursia settled at Monte Cassino
(Barmby, Prolegomena). By then the Emperor Justinian's long campaign to
recover Italy from the Ostrogoths (535--554) had turned the peninsula into a
wasteland of sieges, ruined aqueducts, broken estates, and intermittent
pestilence. Rome itself changed hands repeatedly during Gregory's childhood.
When imperial restoration finally arrived, it brought heavy taxation from
Constantinople rather than the old western prosperity, and in 568---while
Gregory was in his late twenties---the Lombards under King Alboin crossed the
Alps.

Within a generation the Lombards controlled most of northern Italy from their
royal seat at Ticinum (Pavia), along with the powerful southern duchies of
Spoleto and Benevento. Imperial authority shrank to the Exarchate of Ravenna
on the Adriatic, a coastal corridor linking Ravenna to Rome, Naples, southern
Calabria and Apulia, and the islands of Sicily, Sardinia, and Corsica. Between
imperial enclaves and Lombard territories lay contested farmland where no
peasant could count on a safe harvest. Gregory later calculated the Lombard
presence in Italy as a standing era of affliction---twenty-seven years by
594--595 and thirty-five years by 602--603 (Barmby, Prolegomena, citing
\work{Ep.}\ V.21 and XIII.38)---and lived his whole pontificate ``among the
swords of the Lombards.''

\subsection{Family, lineage, and holy women}

Into that fragile world Gregory was born of old Roman wealth and deep
ecclesiastical roots. His contemporary Gregory of Tours, writing from the
eyewitness report of a deacon who visited Rome in 590, states plainly that
Gregory ``descended from one of the first among senatorial families, and lived
in the fear of God from his youth'' (\work{Historia Francorum} X.1). For the
inner life of his household, however, the sharpest witness is Gregory himself.
Preaching in the Roman basilica of Saint Clement around 591--592 on the parable
of the wedding feast---``many are called, but few are chosen''---he paused to
tell his congregation a story from his own family, ``not things done long ago,
but matters for which witnesses still live, who report that they were present''
(\work{Homiliae in Evangelia} 38.15):

\begin{quote}
My father had three sisters, all three of whom were consecrated virgins: one
was called Tharsilla, another Gordiana, and the third Aemiliana. Converted with
one ardor and consecrated at one and the same time, living under regular
discipline, they led a common life in their own house. (\work{Homiliae in
Evangelia} 38.15)
\end{quote}

Gregory goes on to recount a detail that illuminates both his family's piety
and his ancestral connection to the papacy. Tharsilla, distinguished among her
sisters by unceasing prayer and austere abstinence, saw in a night vision
``my forefather Felix, bishop of this Roman Church'' (\latin{atavus meus Felix
hujus Romanae Ecclesiae antistes})---Pope Felix III, who ruled Rome from 483 to
492---beckoning her into a dwelling of everlasting light. Struck soon after
with fever, she lay dying surrounded by kinsfolk, ``among whom my mother also
was present.'' Looking upward, Tharsilla cried out, ``Stand back, stand back,
Jesus is coming!'' (\latin{Recedite, recedite, Jesus venit}). When the women
prepared her body for burial, they found the skin of her knees and elbows
hardened like a camel's from long kneeling in prayer. Before the feast of
Epiphany, Tharsilla appeared to Aemiliana and summoned her too to God, warning
sadly that their third sister, Gordiana, had grown lukewarm and ``been numbered
among laywomen.'' Left alone, Gordiana abandoned her vow of consecration and
married the steward who farmed her lands (\work{Homiliae in Evangelia} 38.15).
Gregory told the story of his own aunt's fall without embarrassment, as a
sobering lesson against spiritual complacency.

Later Roman tradition, preserved by Gregory's ninth-century biographer John the
Deacon and reported in Barmby's Prolegomena, supplies the names of Gregory's
parents: his father Gordianus held the title \latin{regionarius} (an office
linked to ecclesiastical or civic administration in Rome's regions), and his
mother Silvia retired after her husband's death to a solitary cell near the
Aventine, from which she sent raw vegetables to her son in his monastery. Both
Silvia and Gregory's aunts Tharsilla and Aemiliana entered the Roman calendar
of saints.

\witnesslimit{Later medieval and modern accounts routinely assign Gregory to
the illustrious senatorial \latin{gens Anicia}, the aristocratic clan of
Boethius and Saint Benedict's own reputed kin. Neither Gregory's surviving
writings nor Gregory of Tours names the \latin{gens Anicia}. Gregory explicitly
identifies Pope Felix III as his \latin{atavus} (\work{Homiliae in Evangelia}
38.15), and Gregory of Tours confirms his high senatorial birth (\work{Historia
Francorum} X.1); the specific Anician affiliation rests on later tradition and
prosopographical inference rather than direct contemporary attestation.}

\subsection{Education and civic magistracy}

Gregory received the standard rhetorical and legal education of a late Roman
aristocrat. Gregory of Tours reports that ``so accomplished was he in grammar,
dialectic, and rhetoric, that he was held second to none in all the city''
(\work{Historia Francorum} X.1). Legal training gave him the precise grasp of
Roman property law, contracts, wills, prescription periods, and imperial
rescripts that marks every book of his later correspondence. Yet sixth-century
Roman education had narrowed dramatically since the days of Jerome or
Boethius. Gregory never learned Greek. Even after spending six years at the
imperial court in Constantinople, he had to rely on Latin translations and
Eastern correspondents to check Greek patristic authorities---admitting candidly
to Patriarch Eulogius of Alexandria that he knew nothing of the fourth-century
heretic Eudoxius ``in the Latin language'' until Eulogius supplied testimonies
from Basil, Gregory, and Epiphanius (\work{Ep.}\ VII.34, VIII.30; cf.\ Barmby,
Prolegomena, citing \work{Ep.}\ X.35, 39).

Nor did he cherish pagan belles-lettres for their own sake. In the dedicatory
epistle to his \work{Moralia in Iob} he disclaimed the ornaments of classical
rhetoric and declared it unworthy ``to submit the words of the divine Oracle to
the rules of Donatus'' (\work{Moralia}, Ep.\ Ded.\ 5), and as pope he rebuked
Bishop Desiderius of Vienne for lecturing on grammar and pagan literature
(reported by Barmby from \work{Ep.}\ XI.54). His mind was steeped instead in
the Latin Scriptures---using both Jerome's Vulgate and the older Latin
versions---and in the Latin Fathers, above all Saint Augustine, alongside
Ambrose, Jerome, Leo the Great, and Cassian.

Still in his early thirties, certainly by 573, Gregory rose under Emperor
Justin II to the dignity of Prefect of the City of Rome (\latin{Praefectus
Urbi}, styled \latin{Praetor Urbanus} in Barmby's summary). Even in a reduced
city, the urban prefecture remained the highest civil magistracy in Rome,
carrying responsibility for public order, grain supply, civic courts, and
ceremonial precedence in the Senate. Gregory of Tours captures the visual
splendor of those secular years in one memorable sentence: Gregory had been
accustomed ``to pass through the city in the silken garb of high office
adorned with glittering gems'' (\work{Historia Francorum} X.1). Outwardly he
walked in the silk and purple-striped \latin{trabea} of a Roman magistrate;
inwardly, as he later confessed, a quiet struggle had already begun.
